As mentioned in introduction section, in this paper, we formulate maximum likelihood estimation as a convex optimization problem with $l_1$ penalty using Generalized Method of Moments. In this section we describe some existing similar studies as follows. 
In the field of Generalized Method of Moments for estimation we can refer to \cite{lewbel2007local} in which a local Generalized Method of Moments Estimator is proposed for non-parametrically estimation of unknown functions that are defined by conditional moment restrictions. But this research is not about the density estimate and it just finds some function having moment constraints. In \cite{gopalakrishnaneecs} a method with the ability of generating kernel density estimates with a specified degree of sparsity is developed. The proposed method uses $l_0$, negative $l_2$ and weighted $l_1$ penalties in order to explicitly bring sparsity constraint into the optimization problem of estimating density. The author's state that this method can provide sparse solutions while satisfying constraints on the quality of the estimates and therefore provides a trade-off between the sparsity and the quality of the density estimates. The objective function to be minimized in this method is as
\begin{equation}
\label{e1}
\min_a{\alpha^TQ\alpha-c^T\alpha}
\end{equation}
Where,  $\alpha$ is the vector of weights. $Q_ij=k_{\sqrt{2\sigma}}(x_i,x_j), c_i=(2/m)\sum_{j=1}^mk_{\sigma}(y_j,x_i), x_i$-s are independent realization of \emph{f} (density to be estimated) used as training data and $y_i$-s are independent realization of f used as test data. This method uses data samples in the optimization formulation, while our method uses a low dimensional representation of data in the estimation process.
In \cite{hong2013sparse} a sparse kernel density estimator using a forward constrained regression framework is proposed, within which the nonnegative and summing-to-unity constraints of the mixing weights can easily be satisfied. The main contribution of this work is to derive a recursive algorithm to select significant kernel so neat time based on the minimum integrated square error (MISE) criterion for both these lection of kernels and the estimation of mixing weights. The complexity of the proposed algorithm is in the order of the number of training data $N$, which is much lower than the order of $N^2$ offered by the best existing sparse kernel density estimators. This method adaptively fits a linear combination of (gaussian) kernels into the sample data itself. In \cite{peherstorfer2014density} an adaptive sparse-grid-based density estimation method is presented which discretizes the estimated density function on basis functions centered at grid points rather than on kernels centered at the data points. Thus, the costs of evaluating the estimated density function are independent from the number of data points. This method also uses a cross validation technique for the parameter selection. Van Kerm in \cite{van2003adaptive}, proposes an adaptive kernel density estimation using an iterative procedure: An initial (fixed bandwidth) density estimate is computed to get an idea of the density at each of the data points, and this pilot estimate is used to adapt the size of the bandwidth over the data points when computing a new kernel density estimate. In \cite{botev2010kernel} a new adaptive kernel density estimator based on linear diffusion processes is proposed which builds on existing ideas for adaptive smoothing by incorporating information from a pilot density estimate. In addition, a new plug-in bandwidth selection method is proposed that is free from the arbitrary normal reference rules used by existing methods. In \cite{gine2003l1}, the notion of an l1-norm density estimator process indexed by a class of kernels is introduced. Then a functional central limit theorem and a Glivenko–Cantelli theorem are established for this process. While assembling the necessary machinery to prove these results, a body of Poissonization techniques and restricted chaining methods is developed which is useful for studying weak convergence of general processes indexed by a class of functions. None of the theorems imposes any condition at all on the underlying Lebesgue density f. Also, somewhat unexpectedly, the distribution of the limiting Gaussian process does not depend on f. \cite{bunea2010spades} Studies sparse density estimation via $l_1$ penalization (SPADES) with the focus on estimation in high-dimensional mixture models and nonparametric adaptive density estimation. This paper shows that SPADES can recover, with high probability, the unknown components of a mixture of probability densities and that it yields minimax adaptive density estimates. These results are based on a general sparsity oracle inequality that the SPADES estimates satisfy. We offer a data driven method for the choice of the tuning parameter used in the construction of SPADES. The method uses the generalized bisection method introduced in \cite{bunea2009dimension}. The suggested procedure bypasses the need for a grid search and offers substantial computational savings. In \cite{mahapatruni2011cake} a generalization of kernel density estimation called Convex Adaptive Kernel Density Estimation (CAKE), with the aim of minimizing the $l_2$ error of an estimator also known as the integrated squared error (ISE), is presented that replaces single bandwidth selection by a convex aggregation of kernels at all scales, where the convex aggregation is allowed to vary from one training point to another.
A sparse reduced set density estimator with weighted $l_1$ penalty term (RSDE-WL1) is proposed in \cite{yang2015sparser}. By introducing the weighted $l_1$ norm of the estimated coefficients as additional penalty term, the sparse weights of density estimation are estimated, in which small weight coefficients are more likely to be driven to zero. The proposed iterative algorithm is used to solve the corresponding convex optimization problem efficiently. Some discussions about the choice of parameters and the trade-off between sparsity and accuracy are also given. In \cite{bunea2007sparse} it is shown that the penalized least squares estimator satisfies sparsity oracle inequalities, i.e., bounds in terms of the number of non-zero components of the oracle vector. This paper considers estimators of density functions f that belong to the linear span of a set of $M$ basis functions $\{f_1, \ldots, f_M\}$. Let $f_{\lambda}$  denotes the linear combinations:

\begin{equation}
\label{e2}
f_{\lambda}(x)=\sum_{j=1}^M \lambda_jf_j  \hspace{2cm} \lambda=(\lambda_1,\ldots, \lambda_M) \in R^M
\end{equation}
  
This paper defines  $\hat{\lambda}$  ¸ as a minimizer of an $l_1$-penalized criterion, that is called SPADES (SPArse Density Estimation). \par
In \cite{hong2016sparse} a new sparse kernel density estimator with tunable kernels is introduced within a forward constrained regression framework where, by the nonnegative and summing-to-unity constraints of the mixing weights can easily be satisfied. Based on the minimum integrated square error criterion, a recursive algorithm is developed to select significant kernels one at time, and the kernel width of the selected kernel is then tuned using the gradient descent algorithm. In \cite{hong2010sparse} a sparse kernel density estimator is derived based on the zero-norm constraint, in which the zero-norm of the kerne1 weights is incorporated to enhance model sparsity. The classical Parzen window estimate is adopted as the desired response for density estimation, and an approximate function of the zero-norm is used for achieving mathematical tractability and algorithmic efficiency. Under the mild condition of the positive definite design matrix, the kernel weights of the proposed density estimator based on the zero-norm approximation can be obtained using the multiplicative nonnegative quadratic programming algorithm. In \cite{krauthausen2011sparse} the estimation of conditional densities of continuous random variables from noisy samples is considered. The conditional densities are modeled as heteroscedastic Gaussian mixture densities allowing for closed-form solution of Bayesian inference with full densities. The key idea is a regularization based on the curvature of the conditional density function’s surface. The main contributions are the introduction of a superficial regularizer, the consideration of model uncertainty relative to the local data distribution by means of adaptive covariances, and an efficient distance-based estimation algorithm leading to an improved generalization quality of the estimates. The proposed algorithm is an iterative two-step optimization scheme for hyperparameters and the components’ parameters. The obtained solutions are sparse, smooth, and generalize well.
None of above studies has the convex optimization and sparsity and adaptation of bandwidth together in a convex optimization framework on a low dimensional representation of the data. The most similar work to us is \cite{hong2016sparse} which we compare our results to this method in simulation section of this paper.